\documentclass[12pt]{article}
\usepackage[utf8]{inputenc}
\usepackage[T1]{fontenc}
\usepackage{hyperref}
\title{Youjiang Zhuang Language Resources for K-12 Curriculum Development}
\author{Compiled by Claude Sonnet 5.5 (claude-sonnet-5-5) for Andr\'as Kornai}
\date{}
\begin{document}
\maketitle

\section{Introduction}
Youjiang Zhuang (ISO 639-3 zyj, Glottocode youj1238) is a Northern Tai language of the Kra-Dai family. It is spoken in Guangxi, China, in Tiandong County, Tianyang District and parts of Youjiang District in Baise. Wikipedia gives about 870,000 speakers (2007 data) and says speakers call it Gangjdoj, ``local language''. The search found no school curricula or textbooks specific to this language, so little K-12 material is documented here.

\section{Primary Educational Resources}
No primers, textbooks or school programmes specific to Youjiang Zhuang turned up in the search results. This section is therefore empty.

\section{Linguistic and Reference Materials}
\begin{itemize}
\item Wikipedia article on Youjiang Zhuang, which mentions Chinese characters and a romanization as writing systems (one claim is marked ``citation needed''): \url{https://en.wikipedia.org/wiki/Youjiang_Zhuang}
\item Wikitongues language page: \url{https://wikitongues.org/languages/zyj/}
\end{itemize}

\section{Religious and Community Resources}
Global Recordings Network lists Bible-based audio for the related Baise and Duan Xian varieties. It lists no recordings for the language as a whole.
\begin{itemize}
\item Zhuang: Baise (related variety): \url{https://globalrecordings.net/en/language/6034}
\item Zhuang: Duan Xian (related variety): \url{https://globalrecordings.net/en/language/6035}
\end{itemize}

\section{Digital and Audio Resources}
\begin{itemize}
\item Global Recordings Network page, ``Zhuang, Youjiang'': \url{https://www.globalrecordings.net/en/language/zyj}
\end{itemize}

\section{Curriculum Development Context}
Wikitongues rates the language EGIDS 6a (vigorous). The search found no teaching materials, so a curriculum would probably need to start from the Standard Zhuang romanization or from new community work. This report did not verify either option.

\end{document}